% !PRINTED IX
\weberpage{IX}
\inhaltpar{50}{Realitätsbedingungen}{193}
\inhaltpar{51}{Endliche Gruppen linearer gebrochener Substitutionen. Pole der Gruppen}{195}
\inhaltpar{52}{Die verschiedenen Arten möglicher Gruppen}{199}
\inhaltpar{53}{Transformation der Substitutionen von $G$ auf einfache Formen}{203}
\inhaltpar{54}{Die Grundformen}{205}

\inhaltabschnitt{Achter Abschnitt.}{Die Polyëdergruppen.}
\inhaltpar{55}{Die cyklischen Gruppen und die Diëdergruppen}{209}
\inhaltpar{56}{Die Tetraëdergruppe}{212}
\inhaltpar{57}{Die Octaëdergruppe}{216}
\inhaltpar{58}{Die Ikosaëdergruppe}{220}
\inhaltpar{59}{Die Theiler der Ikosaëdergruppe}{227}
\inhaltpar{60}{Die Grundformen der Ikosaëdergruppe}{230}
\inhaltpar{61}{Die Invarianten des Ikosaëders}{232}
\inhaltpar{62}{Polyëdergruppen der zweiten Art. Krystallographische Gruppen}{234}

\inhaltabschnitt{Neunter Abschnitt.}{Congruenzgruppen.}
\inhaltpar{63}{Functionen-Congruenzen}{242}
\inhaltpar{64}{Congruenzkörper}{245}
\inhaltpar{65}{Congruenzgruppen im Körper $\mathfrak{E}$}{250}
\inhaltpar{66}{Einfachheit der Gruppe $E$}{254}
\inhaltpar{67}{Congruenzkörper zweiten Grades}{259}
\inhaltpar{68}{Die reelle lineare Congruenzgruppe $L_p$}{261}
\inhaltpar{69}{Imaginäre Form der Gruppe $L_p$}{266}
\inhaltpar{70}{Divisoren der Gruppe $L_p$, deren Grad durch $p$ theilbar ist}{271}
\inhaltpar{71}{Divisoren der Gruppe $L_p$, deren Grad nicht durch $p$ theilbar ist}{273}
\inhaltpar{72}{Constitution der Gruppe $L_7$ vom Grade 168}{282}

\inhaltbuch{Drittes Buch.}{Anwendungen der Gruppentheorie.}

\inhaltabschnitt{Zehnter Abschnitt.}{Allgemeine Theorie der metacyklischen Gleichungen.}
\inhaltpar{73}{Die Resolventen der Compositionsreihen}{289}
\inhaltpar{74}{Metacyklische Gleichungen}{292}
\inhaltpar{75}{Metacyklische Gleichungen, deren Grad eine Primzahlpotenz ist}{297}
\inhaltpar{76}{Darstellung der \gesperrt{Abel}'schen Gruppe $Q$}{298}
\inhaltpar{77}{Analytische Darstellung der Permutationen}{299}
\inhaltpar{78}{Darstellung der metacyklischen Gruppe $P$}{301}
\inhaltpar{79}{Ternäre lineare Congruenzgruppe für den Modul 2}{306}
\inhaltpar{80}{Resolventen der Gleichung 8\textsuperscript{ten} Grades}{308}
